In this section we present the proposed posterior sampling method. We work at the level of the general affine Gaussian path of Section~\ref{subsec:flow_models}, so that the same construction applies to stochastic interpolants and flow matching, and yields both stochastic (SDE) and deterministic (ODE) samplers from one framework. Throughout, the prior model is assumed trained, providing a velocity $\fmvel_\tau$ (equivalently a score $\genscore_\tau$) for the path that transports a tractable source to the forecast prior $\conddist{\bx_1}{\bx_0}$. The aim is to sample from the one-step posterior $\conddist{\bx_1}{\obs,\bx_0}$ of Eq.~\eqref{eq:posterior_distribution} without retraining.

\subsection{Conditioning the affine Gaussian path on observations}
\label{subsec:conditioning}

The starting point is the observation that conditioning on $\obs$ leaves the affine Gaussian structure of the path intact and only retargets it. Let $\obs = \obsoperator(\bx_1) + \obsnoise$ and write the likelihood $\conddist{\obs}{\bx_1}$. Define the \emph{likelihood tilt}
\begin{align}\label{eq:likelihood_tilt}
    \Phi_\tau^{\obs}(\bx,\bx_0) := \conddist{\obs}{\bX_\tau=\bx,\bx_0}
    = \E\!\left[\conddist{\obs}{\bx_1}\,\middle|\,\bX_\tau=\bx,\bX_0=\bx_0\right],
\end{align}
the probability of the observation given the current state, where the expectation is over the path coupling $\conddist{\bx_1}{\bx_\tau,\bx_0}$. By Bayes' rule the conditional (posterior) marginals factor as
\begin{align}\label{eq:posterior_marginal}
    p_\tau^{\obs}(\bx\mid\bx_0) := \conddist{\bx_\tau=\bx}{\obs,\bx_0}
    = \frac{\Phi_\tau^{\obs}(\bx,\bx_0)\,p_\tau(\bx\mid\bx_0)}{\conddist{\obs}{\bx_0}},
\end{align}
so the posterior path is the prior path reweighted by the likelihood tilt. The next proposition records that this reweighted path is again an affine Gaussian path with the \emph{same} schedules, only with the target replaced by the posterior $\conddist{\bx_1}{\obs,\bx_0}$. All proofs of this section are in Appendices~\ref{appendix:proof_conditioning}--\ref{appendix:proof_multiplicative}.

\begin{proposition}[Posterior path]\label{prop:conditional_path}
Conditioning the interpolant \eqref{eq:general_interpolant} on $\obs$ yields the affine Gaussian path
\begin{align}
    \bx_\tau = \alpha_\tau\bm a_0 + \sigpath_\tau\latentnoise + \beta_\tau\bx_1,
    \qquad \bx_1\sim\conddist{\bx_1}{\obs,\bx_0},\ \ \latentnoise\sim\mathcal{N}(0,\bI),
\end{align}
with $\latentnoise$ still standard normal and independent of $\bx_1$. Consequently its score and velocity are the prior score and velocity shifted by the likelihood score,
\begin{align}
    \genscore_\tau^{\obs}(\bx,\bx_0)
    &= \genscore_\tau(\bx,\bx_0) + \nabla_{\bx}\log\Phi_\tau^{\obs}(\bx,\bx_0),
       \label{eq:conditional_score}\\
    \fmvel_\tau^{\obs}(\bx,\bx_0)
    &= \fmvel_\tau(\bx,\bx_0) + \vscoef_\tau\,\nabla_{\bx}\log\Phi_\tau^{\obs}(\bx,\bx_0),
       \label{eq:conditional_velocity}
\end{align}
where $\vscoef_\tau$ is the velocity--score coefficient \eqref{eq:vscoef_def}. The velocity correction is the score correction scaled by $\vscoef_\tau$.
\end{proposition}

The score relation \eqref{eq:conditional_score} is the standard Bayesian decomposition used in score-based conditional sampling \cite{song_score-based_2021}; the velocity relation \eqref{eq:conditional_velocity} follows from it by applying the velocity--score duality \eqref{eq:general_velocity_of_score} to the posterior path. Because the affine part of \eqref{eq:general_velocity_of_score} depends only on the schedules and the deterministic anchor, which are unchanged by conditioning, it cancels in the difference $\fmvel_\tau^{\obs}-\fmvel_\tau$, leaving the single term $\vscoef_\tau\nabla_{\bx}\log\Phi_\tau^{\obs}$. Equations~\eqref{eq:conditional_score}--\eqref{eq:conditional_velocity} are the only objects the samplers below need.

\subsection{A unified family of posterior samplers}
\label{subsec:posterior_dynamics}

The posterior path has its own marginal-preserving SDE family, obtained by applying Proposition~\ref{prop:marginal_preservation} to $p_\tau^{\obs}$ in place of $p_\tau$. This is the unified posterior sampler.

\begin{theorem}[Unified posterior sampler]\label{theorem:unified_posterior}
Let the prior model define an affine Gaussian path with velocity $\fmvel_\tau$ and score $\genscore_\tau$ whose terminal law is $\conddist{\bx_1}{\bx_0}$, and let $\Phi_\tau^{\obs}$ be the likelihood tilt \eqref{eq:likelihood_tilt}. For any diffusion schedule $\gdiff_\tau\ge0$, the SDE
\begin{align}\label{eq:unified_posterior_sde}
    \rd\bXstar_\tau
    = \left[\fmvel_\tau^{\obs}(\bXstar_\tau,\bx_0) + \tfrac12\gdiff_\tau^2\,\genscore_\tau^{\obs}(\bXstar_\tau,\bx_0)\right]\rd\tau
    + \gdiff_\tau\,\rd\bW_\tau,
    \qquad \bXstar_0\sim p_0^{\obs},
\end{align}
initialised from the reweighted source $p_0^{\obs}(\rd\bx)\propto\Phi_0^{\obs}(\bx,\bx_0)\,p_0(\rd\bx)$, has marginals $p_\tau^{\obs}$ for all $\tau\in[0,1]$, and in particular generates terminal states $\bXstar_1\sim\conddist{\bx_1}{\obs,\bx_0}$.
\end{theorem}

Substituting the conditional velocity and score \eqref{eq:conditional_score}--\eqref{eq:conditional_velocity} into \eqref{eq:unified_posterior_sde} separates the dynamics into the prior drift and a likelihood-score correction,
\begin{align}\label{eq:unified_posterior_drift}
    \rd\bXstar_\tau
    = \underbrace{\left[\fmvel_\tau(\bXstar_\tau,\bx_0) + \tfrac12\gdiff_\tau^2\,\genscore_\tau(\bXstar_\tau,\bx_0)\right]}_{\text{prior drift }\drift^{\gdiff}_\tau}\rd\tau
    + \underbrace{\left(\vscoef_\tau + \tfrac12\gdiff_\tau^2\right)}_{=:\,\gweight_\tau}\nabla_{\bXstar_\tau}\log\Phi_\tau^{\obs}(\bXstar_\tau,\bx_0)\,\rd\tau
    + \gdiff_\tau\,\rd\bW_\tau .
\end{align}
The correction is a diffusion-scaled likelihood score that pushes the process towards regions of high observation likelihood, while the prior drift pulls towards the prior; together they steer the marginals from $p_0^{\obs}$ to the posterior. The single \emph{guidance weight}
\begin{align}\label{eq:guidance_weight}
    \gweight_\tau = \vscoef_\tau + \tfrac12\gdiff_\tau^2
\end{align}
controls the strength of the correction and is the only quantity that changes between the three samplers below. The deterministic part of the weight, $\vscoef_\tau$, comes from the velocity--score duality of the path; the stochastic part, $\tfrac12\gdiff_\tau^2$, comes from the chosen diffusion. Setting $\gdiff_\tau=0$ retains the duality term and yields a deterministic posterior sampler -- the guided probability-flow ODE -- whereas $\gdiff_\tau>0$ adds noise.

\begin{remark}[Relation to Doob's $h$-transform]\label{rem:doob}
An alternative route to a posterior \emph{SDE} fixes the prior diffusion $\gdiff_\tau>0$ and applies Doob's $h$-transform with $h(\bx)=\conddist{\obs}{\bx}$ to its paths, giving the correction $\gdiff_\tau^2\nabla\log\Phi_\tau^{\obs}$ (Theorem~\ref{theorem:posterior_sde}, Appendix~\ref{appendix:proof_posterior}). This path-space construction matches the posterior terminal law for any $\gdiff_\tau>0$, but conditions trajectories rather than marginals, so its intermediate marginals are not $p_\tau^{\obs}$, and it has no deterministic ($\gdiff_\tau=0$) member. The marginal construction \eqref{eq:unified_posterior_drift} instead reproduces $p_\tau^{\obs}$ at every $\tau$ and reduces to the same terminal law; crucially, its weight $\gweight_\tau$ is the one consistent with the tractable interpolated-observation likelihood introduced in Section~\ref{subsec:obs_interpolation} (see Appendix~\ref{appendix:proof_posterior}). We therefore adopt \eqref{eq:unified_posterior_drift} as the working sampler.
\end{remark}

\subsection{Three samplers: stochastic interpolants, flow matching, and the guided ODE}
\label{subsec:supplying_diffusion}

The family \eqref{eq:unified_posterior_drift} contains three samplers of practical interest, differing only in the diffusion $\gdiff_\tau$ (hence the weight $\gweight_\tau$) and in how the velocity and score are read off the trained model.

\paragraph{Sampler 1: stochastic interpolant SDE (native diffusion).}
The SI generator is natively the SDE \eqref{eq:SI_SDE} with diffusion $\gamma_\tau$, whose drift already equals the prior drift $\drift^{\gdiff}_\tau$ with $\gdiff_\tau=\gamma_\tau$, since $\drift_\theta=\fmvel_\tau+\tfrac12\gamma_\tau^2\genscore_\tau$. Choosing $\gdiff_\tau=\gamma_\tau$, Eq.~\eqref{eq:unified_posterior_drift} becomes
\begin{align}\label{eq:posterior_sde_si}
    \rd\bXstar_\tau
    = \drift_\theta(\bXstar_\tau,\bx_0,\tau)\,\rd\tau
    + \gweight_\tau^{\mathrm{SI}}\,\nabla_{\bXstar_\tau}\log\Phi_\tau^{\obs}(\bXstar_\tau,\bx_0)\,\rd\tau
    + \gamma_\tau\,\rd\bW_\tau,
    \qquad
    \gweight_\tau^{\mathrm{SI}} = \vscoef_\tau + \tfrac12\gamma_\tau^2,
\end{align}
with the SI source being the point mass $p_0=\delta_{\bx_0}$, for which the reweighting in Theorem~\ref{theorem:unified_posterior} is vacuous and $\bXstar_0=\bx_0$. No score recovery is needed because the diffusion, and hence the prior drift, is built into the model.

\paragraph{Sampler 2: flow matching SDE (lifted diffusion).}
The FM generator is the deterministic ODE \eqref{eq:fm_ode}, the member $\gdiff_\tau=0$ of the prior family. We lift it to an SDE by selecting any $\gdiff_\tau>0$ and recovering the score from the trained velocity through Eq.~\eqref{eq:fm_score}:
\begin{align}\label{eq:posterior_sde_fm}
    \rd\bXstar_\tau
    = \left[\fmvel_\theta(\bXstar_\tau,\state^{n-1},\tau) + \tfrac12\gdiff_\tau^2\genscore_\tau(\bXstar_\tau)\right]\rd\tau
    + \gweight_\tau^{\mathrm{FM\text-SDE}}\,\nabla_{\bXstar_\tau}\log\Phi_\tau^{\obs}(\bXstar_\tau)\,\rd\tau
    + \gdiff_\tau\,\rd\bW_\tau,
\end{align}
with $\gweight_\tau^{\mathrm{FM\text-SDE}}=\vscoef_\tau+\tfrac12\gdiff_\tau^2$. The diffusion $\gdiff_\tau$ is a free hyperparameter; a schedule vanishing at the endpoints (e.g.\ $\gdiff_\tau\propto\sqrt{\alpha_\tau\beta_\tau}$) keeps the boundary behaviour analogous to SI. Such an endpoint-vanishing choice is also what keeps the lifted drift finite as $\tau\to1$: the recovered FM score \eqref{eq:fm_score} diverges there (its denominator carries $\alpha_\tau\to0$), and only a $\gdiff_\tau$ vanishing fast enough makes the score term $\tfrac12\gdiff_\tau^2\genscore_\tau$ bounded (for rectified flow $\gdiff_\tau^2\propto\alpha_\tau\beta_\tau$ exactly cancels the $1/(1-\tau)$ factor). A constant or slowly-vanishing $\gdiff_\tau$ would expose this singularity. The FM source is $p_0=\mathcal{N}(0,\bI)$. For the marginal construction this initialisation is \emph{exact}, with no reweighting required: at $\tau=0$ we have $\beta_0=0$ and $\alpha_0=1$, so $\bx_0=\latent$ is independent of $\bx_1$ under the interpolant coupling, the tilt $\Phi_0^{\obs}(\bx)=\E[\conddist{\obs}{\bx_1}\mid\bx_0]$ is constant in $\bx$, and hence $p_0^{\obs}\propto\Phi_0^{\obs}\,p_0=\mathcal{N}(0,\bI)$ (Appendix~\ref{appendix:proof_conditioning}). The reweighting becomes non-trivial only in the path-space Doob construction (Theorem~\ref{theorem:posterior_sde}), where the Markov SDE correlates $\bXstar_0$ with $\bXstar_1$ and $\Phi_0^{\obs}$ is genuinely state-dependent.

\paragraph{Sampler 3: guided probability-flow ODE (deterministic).}
Setting $\gdiff_\tau=0$ in \eqref{eq:unified_posterior_drift} gives a deterministic posterior sampler driven purely by the velocity--score duality term,
\begin{align}\label{eq:posterior_ode_fm}
    \rd\bXstar_\tau
    = \fmvel_\theta(\bXstar_\tau,\state^{n-1},\tau)\,\rd\tau
    + \vscoef_\tau\,\nabla_{\bXstar_\tau}\log\Phi_\tau^{\obs}(\bXstar_\tau)\,\rd\tau,
    \qquad \gweight_\tau^{\mathrm{ODE}} = \vscoef_\tau,
\end{align}
which is exactly the probability-flow ODE of the posterior path \eqref{eq:posterior_marginal}: its velocity is the conditional velocity \eqref{eq:conditional_velocity}. The guidance term coincides with the interpolant-guided flow used for linear inverse problems \cite{yan_fig_2024, pokle_training-free_2024}, here derived as the $\gdiff_\tau=0$ member of the same family and specialised to data assimilation. Because the dynamics are deterministic, all randomness must enter through the source: the ODE transports the source to the posterior. For flow matching this source is simply $p_0=\mathcal{N}(0,\bI)$, which already coincides with $p_0^{\obs}$ since $\Phi_0^{\obs}$ is constant (no initial reweighting is needed); the guidance term bends the otherwise prior-sampling flow towards the posterior along the path. This is well posed for flow matching, whose source $\mathcal{N}(0,\bI)$ is non-degenerate, but \emph{not} for stochastic interpolants, whose point-mass source $\delta_{\bx_0}$ would be mapped deterministically to a single point rather than spread over the posterior. The guided ODE is therefore the flow matching counterpart of the two SDE samplers, sharing their guidance machinery while removing the injected noise.

Table~\ref{tab:samplers} summarises the three samplers. They use identical observation-handling (Sections~\ref{subsec:obs_interpolation}--\ref{sec:multiplicative_correction}) and differ only in the boxed quantities $\gdiff_\tau$, $\gweight_\tau$, and the source.

\subsection{Tractable likelihood via observation interpolation}
\label{subsec:obs_interpolation}

Evaluating the correction in any of \eqref{eq:posterior_sde_si}--\eqref{eq:posterior_ode_fm} requires the likelihood score $\nabla_{\bx}\log\Phi_\tau^{\obs}$ at the intermediate states $\bx_\tau$. The exact tilt,
\begin{align}\label{eq:phi_integral}
    \Phi_{\tau}^{\obs}(\bx_\tau, \bx_0)
    = \int \conddist{\obs}{\bx_1}\, \conddist{\bx_1}{\bx_\tau, \bx_0}\, \rd \bx_1,
\end{align}
is intractable: it requires sampling $\conddist{\bx_1}{\bx_\tau, \bx_0}$ by integrating the dynamics forward in pseudo-time and backpropagating through the trajectory. We avoid this by interpolating the observation itself. Define the \emph{observation interpolant}
\begin{align} \label{eq:interpolated_observation}
    \interpolantobs_\tau(\bx_0, \obs) = \alpha_\tau \obsmatrix\bm{a}_0 + \beta_\tau \obs,
    \qquad \interpolantobs_0 = \obsmatrix\bm a_0,\quad \interpolantobs_1 = \obs,
\end{align}
which mirrors the state interpolant in observation space: it transports the observation operator applied to the deterministic anchor at $\tau=0$ to the true observation at $\tau=1$. For SI ($\bm a_0=\bx_0$) this is $\interpolantobs_\tau = \alpha_\tau\obsmatrix\bx_0 + \beta_\tau\obs$, while for FM ($\bm a_0=\bm 0$) it reduces to $\interpolantobs_\tau = \beta_\tau\obs$. The likelihood tilt at pseudo-time $\tau$ is then replaced by the directly computable surrogate
\begin{align}
    \bar{\Phi}_{\tau}^{\interpolantobs_\tau}(\bx_\tau, \bx_0)
    = \conddist[\tau]{\interpolantobs_\tau}{\bx_\tau, \bx_0},
\end{align}
giving the posterior drift correction in terms of the interpolant-likelihood score $\bar S_\tau := \nabla_{\bx_\tau}\log\bar\Phi_\tau^{\interpolantobs_\tau}$,
\begin{align} \label{eq:interpolant_posterior_drift}
    \postdrift_\theta(\bx_{\tau}, \bx_0, \interpolantobs_\tau, \tau)
    =
    \drift^{\gdiff}_\tau(\bx_{\tau}, \bx_0)
    + \gweight_\tau\,
    \nabla_{\bx_\tau} \log\bar{\Phi}_{\tau}^{\interpolantobs_\tau}(\bx_\tau, \bx_0).
\end{align}
This is the working form of \eqref{eq:unified_posterior_drift}, valid for all three samplers through the corresponding weight $\gweight_\tau$.

To use this surrogate we need the conditional law of $\interpolantobs_\tau$ given $\bx_\tau$. Writing the zero-mean source process along the path as $\src_\tau := \sigpath_\tau\latentnoise$ (so that $\bx_\tau = \alpha_\tau\bm a_0 + \beta_\tau\bx_1 + \src_\tau$), the following holds for both model classes.

\begin{lemma}[Interpolated observation likelihood]\label{lem:interpolated_likelihood}
Let $\obsoperator$ be linear with matrix $\obsmatrix$ and the observation noise be Gaussian, $\obsnoise \sim \mathcal{N}(0,\obscov)$. Then the interpolated observation satisfies
\begin{equation}\label{eq:ybar_decomp}
  \interpolantobs_\tau = \obsmatrix\bx_\tau + \interpolantobsnoise_\tau,
  \qquad
  \interpolantobsnoise_\tau := \beta_\tau\obsnoise - \obsmatrix\src_\tau,
\end{equation}
and the conditional mean and covariance of $\conddist{\interpolantobs_\tau}{\bx_0, \bx_\tau}$ are
\begin{align}
  \E[\interpolantobs_\tau \mid \bx_\tau, \bx_0]
  &= \obsmatrix\bx_\tau - \obsmatrix\,\srcmean_\tau(\bx_\tau,\bx_0) =: \bar{\mu}_\tau(\bx_\tau, \bx_0),
     \label{eq:bias} \\[3pt]
  \cov(\interpolantobs_\tau \mid \bx_\tau, \bx_0)
  &= \beta_\tau^2\obscov + \obsmatrix\,\srccov_\tau(\bx_\tau,\bx_0)\,\obsmatrix^\top =: \bar{\Sigma}_{\tau}(\bx_\tau, \bx_0),
     \label{eq:cov_correction}
\end{align}
where $\srcmean_\tau := \E[\src_\tau\mid\bx_\tau,\bx_0]$ and $\srccov_\tau := \cov(\src_\tau\mid\bx_\tau,\bx_0)$ are the conditional moments of the source process.
\end{lemma}
The proof is in Appendix~\ref{appendix:proof_interpolant_likelihood}. The source moments are themselves computable from the trained model, with the two model classes differing only in what plays the role of $\src_\tau$.

\begin{lemma}[Source conditional moments]\label{lem:source_moments}
The conditional moments of the source process are:
\begin{itemize}
    \item \emph{Stochastic interpolants}, $\src_\tau = \gamma_\tau\bW_\tau$:
    \begin{align}
        \srcmean_\tau = -\gamma_\tau^2\tau A_\tau\!\left(\beta_\tau\drift_\theta(\bx,\bx_0,\tau) - c_\tau(\bx,\bx_0)\right),
        \quad
        \srccov_\tau = \gamma_\tau^2\tau\bI + \gamma_\tau^4\tau^2 A_\tau\!\left(\beta_\tau J_{\drift_\theta} - \dot\beta_\tau\bI\right),
    \end{align}
    with $A_\tau, c_\tau$ as in Eq.~\eqref{eq:SI_score} and $J_{\drift_\theta}=\nabla_{\bx}\drift_\theta$.
    \item \emph{Flow matching}, $\src_\tau = \alpha_\tau\latent$:
    \begin{align}\label{eq:fm_source_moments}
        \srcmean_\tau = -\alpha_\tau^2\,\genscore_\tau(\bx),
        \qquad
        \srccov_\tau = \alpha_\tau^2\bI + \alpha_\tau^4\,\nabla_{\bx}\genscore_\tau(\bx),
    \end{align}
    with $\genscore_\tau$ given in closed form by Eq.~\eqref{eq:fm_score}.
\end{itemize}
\end{lemma}
Both cases follow from Stein's and the second-order Tweedie identities; see Appendix~\ref{appendix:proof_interpolant_likelihood}. In each case the conditional mean of the interpolated observation equals the observation interpolant evaluated at the model's denoiser, $\bar\mu_\tau = \alpha_\tau\obsmatrix\bm a_0 + \beta_\tau\obsmatrix\,\E[\bx_1\mid\bx_\tau,\bx_0]$, so Eq.~\eqref{eq:interpolant_posterior_drift} is a Tweedie-style guidance term shared by all three samplers.

\subsection{Multiplicative correction of the likelihood score}
\label{sec:multiplicative_correction}

The interpolated likelihood is an approximation of the true likelihood tilt $\Phi_\tau^{\obs}$. We quantify the gap in closed form. Assume the interpolant likelihood is Gaussian,
\begin{align}\label{eq:interpolant_likelihood_gaussian}
    \tilde{p}(\interpolantobs_\tau \mid \bx_\tau, \bx_0)
    := \mathcal{N}\!\left(\bar{\mu}_\tau(\bx_\tau, \bx_0),\, \bar{\Sigma}_\tau(\bx_\tau, \bx_0)\right),
\end{align}
and that the true likelihood admits a Gaussian surrogate
\begin{align}\label{eq:true_likelihood_gaussian}
    \Phi_{\tau}^{\obs}(\bx_\tau, \bx_0)
    \approx \mathcal{N}\!\left(\obs;\, \mu_\tau(\bx_\tau, \bx_0),\, \obscov\right),
    \qquad
    \mu_\tau := \obsmatrix\,\E[\bx_1 \mid \bx_\tau, \bx_0],
\end{align}
whose mean is the model denoiser pushed through $\obsmatrix$ and whose covariance is the raw observation noise. This is the Diffusion Posterior Sampling (DPS) surrogate \cite{chung_diffusion_2023, chung_improving_2024, song_score-based_2021}: it approximates the true tilt $\Phi_\tau^{\obs}$ by a Gaussian centred at the denoiser with the \emph{uninflated} covariance $\obscov$. The corrected score $S_\tau$ below is therefore the DPS guidance score, not the genuine non-Gaussian likelihood score. A finer (\(\Pi\)GDM-style) surrogate \cite{song_pseudoinverse_2023} would inflate the covariance to $\obscov + \obsmatrix\,\cov(\bx_1\mid\bx_\tau,\bx_0)\,\obsmatrix^\top$; this term is available at no extra cost here, since $\cov(\bx_1\mid\bx_\tau,\bx_0)=\srccov_\tau/\beta_\tau^2$ is exactly the source covariance of Lemma~\ref{lem:source_moments}. We use the uninflated surrogate for simplicity and leave the inflated variant to future work. Let $\bar{S}_\tau := \nabla_{\bx_\tau} \log \bar{\Phi}_{\tau}^{\interpolantobs_\tau}$ and $S_\tau := \nabla_{\bx_\tau} \log \Phi_{\tau}^{\obs}$ be the interpolant- and surrogate-likelihood scores, where, here and in the following, gradients of the Gaussian log-densities are taken with the covariances $\bar{\Sigma}_\tau$ and $\obscov$ held fixed with respect to $\bx_\tau$, as is common practice in guidance methods.

\begin{theorem}[Multiplicative correction]\label{thm:multiplicative_correction}
Under \eqref{eq:interpolant_likelihood_gaussian}--\eqref{eq:true_likelihood_gaussian},
\begin{align}\label{eq:multiplicative_correction}
    S_\tau = G_\tau(\bx_\tau, \bx_0)\, \bar{S}_\tau,
    \qquad
    G_\tau(\bx_\tau, \bx_0)
    = \bI + \frac{1}{\beta_\tau^2}\,
      \srccov_\tau(\bx_\tau, \bx_0)\,
      \obsmatrix^\top \obscov^{-1} \obsmatrix,
\end{align}
with $\srccov_\tau$ the conditional source covariance of Lemma~\ref{lem:source_moments}. The factor $G_\tau - \bI$ has rank at most $N_y$ and vanishes on $\ker \obsmatrix$.
\end{theorem}

See Appendix~\ref{appendix:proof_multiplicative} for the proof, which is carried out for the general source covariance and so covers all three samplers. Substituting \eqref{eq:multiplicative_correction} into Eq.~\eqref{eq:interpolant_posterior_drift} gives the posterior drift, exact under the Gaussian surrogates \eqref{eq:interpolant_likelihood_gaussian}--\eqref{eq:true_likelihood_gaussian},
\begin{align}\label{eq:posterior_drift_multiplicative}
    \postdrift_\theta(\bx_\tau, \bx_0, \obs, \tau)
    = \drift^{\gdiff}_\tau(\bx_\tau, \bx_0)
    + \gweight_\tau\, G_\tau(\bx_\tau, \bx_0)\, \bar{S}_\tau,
\end{align}
obtained from the interpolant-score drift by a single rank-$N_y$ operator and the scalar weight $\gweight_\tau$.

\begin{corollary}[Jacobian-free posterior drift]\label{cor:cheap_drift}
Dropping the Jacobian term in $\srccov_\tau$ (i.e.\ $\srccov_\tau\approx\gamma_\tau^2\tau\bI$ for SI and $\srccov_\tau\approx\alpha_\tau^2\bI$ for FM) gives the state-independent gain
\begin{align}\label{eq:G_cheap}
    G_\tau \approx \bI + \frac{\rho_\tau}{\beta_\tau^2}\, \obsmatrix^\top \obscov^{-1} \obsmatrix,
    \qquad
    \rho_\tau = \begin{cases} \gamma_\tau^2\tau & \text{(SI)},\\[2pt] \alpha_\tau^2 & \text{(FM)},\end{cases}
\end{align}
which depends on $\tau$ only through the scalar $\rho_\tau/\beta_\tau^2$ and requires no Jacobian. When $\obsmatrix$ and $\obscov$ are time-independent, $\obsmatrix^\top \obscov^{-1} \obsmatrix$ is precomputed once and the correction costs $\mathcal{O}(N_u N_y)$ per step.
\end{corollary}

Dropping the Jacobian replaces $\srccov_\tau$ by its isotropic part and is accurate when $\sigpath_\tau^2\|\nabla_{\bx}^2\log p_\tau\|\ll1$, i.e.\ near $\tau\to1$ (small source scale) or for near-Gaussian priors. For multimodal priors at intermediate $\tau$ the discarded term is the second-order curvature that distinguishes this gain from an isotropic one, so the full $G_\tau$ may be preferable; we treat the choice as an empirical trade-off.

\begin{remark}[Boundary behaviour]
At $\tau\to1$ both source scales vanish ($\gamma_1=0$ for SI, $\alpha_1=0$ for FM) while $\beta_1=1$, so $G_\tau\to\bI$: the interpolant and true scores coincide at the terminal time, consistent with $\interpolantobs_1=\obs$. For $\tau<1$ the gain amplifies $\bar S_\tau$ along directions visible to $\obsmatrix$, with strength set by the source signal-to-noise ratio $\rho_\tau/\beta_\tau^2$ relative to $\obscov$. The multiplicative correction is independent of the sampler: it acts on the likelihood score $\bar S_\tau$, which all three samplers share, and is then scaled by the sampler's weight $\gweight_\tau$.
\end{remark}

\subsection{Summary: one framework, three samplers}
\label{subsec:summary_table}

Tables~\ref{tab:samplers} and~\ref{tab:si_vs_fm} collect the construction. Table~\ref{tab:samplers} contrasts the three samplers, which are the members $\gdiff_\tau=0$ (guided ODE), $\gdiff_\tau=\gamma_\tau$ (SI-SDE), and $\gdiff_\tau>0$ (FM-SDE) of the single family \eqref{eq:unified_posterior_drift}, differing only in the diffusion, the resulting guidance weight $\gweight_\tau=\vscoef_\tau+\tfrac12\gdiff_\tau^2$, and the source. Table~\ref{tab:si_vs_fm} collects the model-specific bookkeeping for the two model classes: what the source process is, how the velocity and score are read off the trained network, and the resulting source moments and gain $G_\tau$. The posterior construction -- conditioned path, observation interpolant, and multiplicative gain -- is identical across samplers.

\begin{table}[ht]
    \centering
    \caption{The three posterior samplers as members of the unified family \eqref{eq:unified_posterior_drift}. All share the prior drift $\drift^{\gdiff}_\tau=\fmvel_\tau+\tfrac12\gdiff_\tau^2\genscore_\tau$, the likelihood-score correction $\gweight_\tau\,G_\tau\bar S_\tau$, and the observation-interpolant machinery; they differ only in the diffusion $\gdiff_\tau$, the guidance weight $\gweight_\tau=\vscoef_\tau+\tfrac12\gdiff_\tau^2$, and the source.}
    \label{tab:samplers}
    \small
    \begin{tabular}{p{0.2\linewidth}|p{0.24\linewidth}|p{0.24\linewidth}|p{0.21\linewidth}}
        \toprule
        \textbf{Property} & \textbf{Guided ODE (FM)} & \textbf{SI-SDE} & \textbf{FM-SDE} \\
        \midrule
        Dynamics & deterministic ODE & stochastic SDE & stochastic SDE \\
        Diffusion $\gdiff_\tau$ & $0$ & native $\gamma_\tau$ & lifted $\gdiff_\tau>0$ \\
        Guidance weight $\gweight_\tau$ & $\vscoef_\tau$ & $\vscoef_\tau+\tfrac12\gamma_\tau^2$ & $\vscoef_\tau+\tfrac12\gdiff_\tau^2$ \\
        Source $p_0$ & $\mathcal{N}(0,\bI)$ & $\delta_{\bx_0}$ & $\mathcal{N}(0,\bI)$ \\
        Reweighted init & not needed ($\Phi_0^{\obs}$ const) & vacuous (point mass) & not needed ($\Phi_0^{\obs}$ const) \\
        Velocity/score & from $\fmvel_\theta$ \eqref{eq:fm_score} & from $\drift_\theta$ \eqref{eq:SI_score} & from $\fmvel_\theta$ \eqref{eq:fm_score} \\
        Score recovery & not needed (ODE) & built into drift & required (lift) \\
        \bottomrule
    \end{tabular}
\end{table}

\begin{table*}[ht]
    \centering
    \caption{Model-specific quantities for the observation-interpolant method, instantiated for stochastic interpolants and flow matching. The general framework (left column) treats both as affine Gaussian probability paths; the model-specific quantities differ only in the source process $\src_\tau$ and how the velocity and score are obtained. The guided ODE and FM-SDE share the flow matching column; the SI-SDE uses the stochastic interpolant column.}
    \label{tab:si_vs_fm}
    \small
    \begin{tabular}{p{0.235\linewidth}|p{0.345\linewidth}|p{0.345\linewidth}}
        \toprule
        \textbf{Component} & \textbf{Stochastic interpolant} & \textbf{Flow matching} \\
        \midrule
        Interpolant $\bx_\tau$
            & $\alpha_\tau\bx_0 + \beta_\tau\bx_1 + \gamma_\tau\bW_\tau$
            & $\alpha_\tau\latent + \beta_\tau\bx_1,\ \ \latent\sim\mathcal{N}(0,\bI)$ \\
        Source $p_0$
            & $\delta_{\bx_0}$ (previous state) $+$ Brownian
            & $\mathcal{N}(0,\bI)$ (latent), $\bx_0$ as input \\
        Source process $\src_\tau$
            & $\gamma_\tau\bW_\tau$
            & $\alpha_\tau\latent$ \\
        Learned object
            & SDE drift $\drift_\theta=\E[\bR_\tau\mid\bx_\tau]$
            & velocity $\fmvel_\theta=\E[\dot\alpha_\tau\latent+\dot\beta_\tau\bx_1\mid\bx_\tau]$ \\
        Score from model
            & $A_\tau[\beta_\tau\drift_\theta - c_\tau]$ \eqref{eq:SI_score}
            & $\dfrac{\beta_\tau\fmvel_\theta-\dot\beta_\tau\bx}{\alpha_\tau(\dot\beta_\tau\alpha_\tau-\dot\alpha_\tau\beta_\tau)}$ \eqref{eq:fm_score} \\
        Velocity--score coef.\ $\vscoef_\tau$
            & $\dfrac{\sigpath_\tau(\dot\beta_\tau\sigpath_\tau-\dot\sigpath_\tau\beta_\tau)}{\beta_\tau},\ \ \sigpath_\tau=\gamma_\tau\sqrt\tau$
            & $\dfrac{\alpha_\tau(\dot\beta_\tau\alpha_\tau-\dot\alpha_\tau\beta_\tau)}{\beta_\tau}$ \\
        Native generator
            & SDE $\rd\bX=\drift_\theta\rd\tau+\gamma_\tau\rd\bW$
            & ODE $\rd\bX=\fmvel_\theta\rd\tau$ \\
        Available samplers
            & SI-SDE ($\gdiff_\tau=\gamma_\tau$)
            & FM-SDE ($\gdiff_\tau>0$), guided ODE ($\gdiff_\tau=0$) \\
        Observation interpolant $\interpolantobs_\tau$
            & $\alpha_\tau\obsmatrix\bx_0 + \beta_\tau\obs$
            & $\beta_\tau\obs$ \\
        Likelihood mean $\bar\mu_\tau$
            & $\obsmatrix\bx_\tau-\gamma_\tau\obsmatrix\,\E[\bW_\tau\mid\bx_\tau]$
            & $\obsmatrix\bx_\tau+\alpha_\tau^2\obsmatrix\genscore_\tau$ \\
        Source moments
            & Wiener moments (Lemma~\ref{lem:source_moments})
            & Tweedie moments via $\genscore_\tau$ \eqref{eq:fm_source_moments} \\
        Multiplicative gain $G_\tau$
            & $\bI+\frac{\gamma_\tau^2}{\beta_\tau^2}\Sigma^{\bW}_\tau\obsmatrix^\top\obscov^{-1}\obsmatrix$
            & $\bI+\frac{\alpha_\tau^2}{\beta_\tau^2}(\bI+\alpha_\tau^2\nabla\genscore_\tau)\obsmatrix^\top\obscov^{-1}\obsmatrix$ \\
        \bottomrule
    \end{tabular}
\end{table*}
